\subsection{Tree building process}

\indent 

In theory, the decision regions could have any shape.

However, we choose to divide the predictor space into high-dimensional rectangles, or boxes, for simplicity and for ease of interpretation of the resulting predictive model.

The goal is to find boxes $R_1,\ldots,R_J$  that minimizes the residual sum of squares (RSS), given by:

\[
\sum_{j=1}^J \sum_{i\in R_j}(y_i - \hat{y}_{R_j})^2
\]

$\hat{y}_{R_j}$ is the mean response for the training observations within the $j^{th}$ box.

\subsubsection{Recursive Binary Splitting}

\indent 

We want to partition the feature space into $J$ distinct regions.

But trying every possible partition is computationally infeasible.

Instead, we use a \textbf{top-down greedy approach }(recursive binary splitting): At each step, choose the best split. Recursively repeat the process on resulting subspaces.

\begin{itemize}
    \item \textbf{Top-down}: Start at the root node (entire dataset) and split down.
    \item \textbf{Greedy}: At each step, choose the split that most reduces RSS or impurity at that step.
    \item \textbf{Recursive}: Once a split is made, repeat the process on each child node.
\end{itemize}

This method does not look ahead — it optimises locally, not globally.

\subsubsection{Overfitting with Recursive Binary Splitting}

\indent 

Trees are grown via recursive binary splitting can lead to deep, complex trees that fit training data perfectly.

Such trees may overfit as it captures noise instead of signal.

This leads to \textbf{high variance} and \textbf{poor generalisation}.

\subsubsection{Tree Pruning}

\indent

Tree pruning is like regularisation — we sacrifice training accuracy to gain prediction stability (lowering the variance).

\noindent\textbf{Cost-Complexity Pruning (the weakest link pruning)}

\[
C_{\alpha}(T) = \sum_{m=1}^{|T|}\sum_{i:x_i\in R_m}(y_i - \hat{y}_{R_m})^2 + \alpha |T|
\]

$T$: subtree. $R_m$: region(leaf) $m$. $\hat{y}_{R_m}$: mean response in region $m$. $\alpha$: tuning parameter controls a trade-off between the subtree’s complexity and its fit to the training data.

\subsection{Using Decision Trees for prediction}

\indent 

We divide the predictor space: That is, the set of possible values for $X_1,\ldots,X_p$, into $J$ distinct and non-overlapping regions, $R_1,\ldots, R_J$.

For every observation that falls into the region $R_j$, we make the same prediction, which is simply the mean of the response values for the training observations in $R_j$.

\subsubsection{Decision Trees for Classification}

\indent 

Very similar to a regression tree. Exception being the prediction of a qualitative response rather than a quantitative one.

For a classification tree, inspect the region that the observation belongs and predict the most commonly occurring class in that region

\subsubsection{Gini index}

\indent 

Just as in the regression setting, we use recursive binary splitting to grow a classification tree. In the classification setting, RSS cannot be used as a criterion for making the binary splits. Alternative measure such as Gini index is used instead. The Gini index is defined by:

\[
G = \sum_{k=1}^K \hat{p}_{mk}(1-\hat{p}_{mk})
\]

where $\hat{p}_{mk}$ represents the proportion of training observations in the $m^{th}$ region that are from the $k^{th}$ class.

The Gini index is a measure of total variance across the $K$ classes. It takes on a small value if all of the $\hat{p}_{mk}$ values are close to zero or one. This occurs when there is a clear majority class!

For this reason the Gini index is referred to as a measure of node \textbf{purity}. A small value indicates that a node contains predominantly observations from a single class.

\subsubsection{Entropy}

\indent

An alternative to the Gini index is entropy, given by:

\[
D = -\sum_{j=1}^J\sum_{k=1}^K \hat{p}_{jk}\log \hat{p}_{jk}
\]

It turns out that the Gini index and the Entropy are very similar numerically.

\subsubsection{Advantages and disadvantages of trees}

\indent

\noindent\textbf{Advantages}

\begin{itemize}
    \item Trees are very easy to explain to people
    \item Some people believe that decision trees closely relate to human decision-making
    \item Trees can be displayed graphically
    \item Can handle different data types and doesn’t require scaling
\end{itemize}

\noindent\textbf{Disadvantages}

Trees do not have the same level of predictive accuracy as some of the other regression and classification approaches we have discussed.

\subsection{Ensemble of trees}

\indent 

Build multiple decision trees. Combine their results to form a more \textbf{stable and accurate} predictor.

\noindent\textbf{Key ideas}:

\begin{itemize}
    \item Take samples of trees (e.g. using bootstrapping)
    \item Use random subsets of features at each split
    \item Aggregate predictions:
    \begin{itemize}
        \item Majority vote (classification)
        \item Average (regression)
    \end{itemize}
\end{itemize}

The strength of ensemble methods lies in \textbf{reducing variance} and \textbf{improving robustness} through aggregation. “Many weak learners can form a strong learner.”

\subsubsection{Bagging (Bootstrap Aggregating)}

\indent 

Bootstrap aggregating, or bagging, is a general purpose procedure for reducing the variance of a statistical learning method. It is particularly useful and frequently used in the context of decision trees.

Recall that given a set of $n$ independent observations $Z_1,\ldots,Z_n$ each with a variance of $\sigma^2$: the variance of the mean $\bar{Z}$ is given by $\sigma^2/n$.

In other words, averaging a set of observations reduces variance.

Since it is typically not possible to have access to multiple training sets. We can use bootstrapping to create multiple training sets.

Use bootstrapping to take repeated samples from a (single) training data set.

In this approach, we generate $B$ different bootstrapped training data sets.

Train our method of the $b^{th}$  bootstrapping set in order to obtain $\hat{f}_b^*(x)$, the prediction at a point $x$.

Average all the observations to obtain:

\[
\hat{f}(x) = \frac{1}{B}\sum_{b=1}^B \hat{f}_b^*(x)
\]

It turns out that there is a very straightforward way to estimate the test error of a bagged model.

Recall that the key to bagging is that trees are repeatedly fit to bootstrapped subsets of the observations. One can show that on average, each bagged tree makes use of around $2/3$ of the observations when the sample size 
$n$ is large.

The remaining $1/3$ of the observations not used to fit a given bagged tree are referred to as the out-of-bag (OOB) observations.

We can predict the response for the $i^{th}$ observation using each of the trees in which that observation was OOB. This will yield around $B/3$ predictions for the $i^{th}$ observation, which we average. This estimate is essentially the leave-one-out (LOO) cross-validation error for bagging, if $B$ is large.

\subsubsection{From bagging to Random Forest}

\indent

Random forests (sometimes) provide an improvement over bagged trees by way of a small tweak that decorrelates the trees. This reduces the variance when we average the trees.

As in bagging, we build a number of decision trees on bootstrapped training samples.

But when building these decision trees, each time a split in a tree is considered, a random selection of $m$ predictors is chosen as split candidates from the full set of $p$ predictors. The split is allowed to use only one of those $m$ predictors.

A fresh selection of $m$ predictors is taken at each split, and typically we choose $m = \sqrt{p}$. That is, the number of predictors considered at each split is approximately equal to the square root of the total number of predictors.

\subsubsection{Boosting}

\indent

Like bagging, boosting is a general approach that can be applied to many statistical learning methods for regression or classification. We only discuss boosting for decision trees.

Recall that bagging involves creating multiple copies of the original training data set using the bootstrap, fitting a separate decision tree to each copy, and then combining all of the trees in order to create a single predictive model.

Notably, each tree is built on a bootstrap data set, independent of other trees.

Boosting works in a similar way, except that the trees are grown sequentially. Each tree is grown using information from previously grown trees.

Unlike fitting a single large decision tree to the data, which amounts to fitting the data hard and potentially overfitting, the boosting approach instead learns slowly. Each new model attempts to \textbf{correct the errors} made by the previous ones. The final model is a \textbf{weighted sum} of weak learners (usually decision trees).

By fitting small trees to the residuals, we slowly improve $\hat{f}(x)$ in areas where it does not perform well.

There is a shrinkage parameter $\lambda$ that slows the process down even further, allowing more and different shaped trees to attack the residuals.

\step $\hat{f}(x) = 0$ and $r_i = y_i$ for all $i$ in the trainging set. $r_i$ denotes the $i^{th}$ residual, and $y_i$ is the outcome.

\step For $b=1,2,\ldots,B$, fit a tree $\hat{f}_b$ with $d$ splits ($d+1$ terminal nodes) to the new training data $(X,r)$. That is, $r$ is the new response value. 

\step Update $\hat{f}$ by adding in a shrunken version of the new tree $\hat{f}(x) \xleftarrow{}\hat{f}(x) + \lambda \hat{f}_b(x)$. 

\step Update the residuals $r_i \xleftarrow{} r_i - \lambda \hat{f}_b(x)$.

\step Compute the final model $\hat{f}(x) = \sum_{i=1}^B \lambda\hat{f}_b(x)$

\bigskip

\noindent\textbf{Parameters to tune in boosting}:

\begin{itemize}
    \item Number of trees $B$
    \item The shrinkage parameter $\lambda$: a small positive number. Typical values are between 0.01 and 0.001.
    \item Number of split $d$ in each tree. Controls the complexity of the boosted ensemble. If $d=1$, then the tree is just a stump. This actually usually works quite well
\end{itemize}

\subsection{Other boosting algorithms}

\subsubsection{AdaBoost}

\indent 

Short for Adaptive Boosting, one of the first boosting algorithms for classification.

Convert a set of weak classifiers (decision stumps) into a strong one

Basic idea: At each iteration, reweight the data to place more weight on data points that the classifier got wrong.

At the end, combine all the weak classifiers by taking a weighted combination. Put more weight on the weak classifiers with the higher accuracies.

\subsubsection{Stochastic gradient boosting}

\indent 

Use the idea behind bagging.

In each iteration of stochastic boosting, a sample of the training set instead of the full training set is used.

Instead of a bootstrap sample (with replacement), the algorithms samples a fraction of the training set.

This introduced randomness can improve performance.

\subsubsection{XGBoost}

\indent 

XGBoost is short for eXtreme Gradient Boosting.

It is a library for high performance gradient boosting models written in C++ with implementations in Python, R, and Julia.

Designed for speed: up to 10 times faster than the gbm package.

Has been very popular in recent years and has won a number of machine learning competitions (especially on tabular data).

Supports regularization.

Can handle missing data.

\subsubsection{Bagging vs Boosting}

\indent 

Both ensemble methods get $N$ learners from 1 learner. Built independently for bagging. Built sequentially for boosting.

Trees built in boosting are weak learners (sometimes just a stump) while trees in Random Forest have higher complexity

Fewer parameters to tune in Random Forest, many more in Boosting (depending on which variation you are using)

Both combine outputs from $N$ trees.

\subsection{Summary}

\indent 

Decision trees are simple and interpretable models for regression and classification.

However, they are often not competitive with other methods in terms of prediction accuracy.

Bagging, random forests, and boosting are good methods for improving the prediction accuracy of trees. They work by growing many trees on the training data and then combining the predictions of the resulting ensemble of trees.

The latter two methods, random forests and boosting, are among the state-of-the-art methods for supervised learning. Their results can be difficult to interpret. But we can still get relative importance of each predictor. For classification, the relative importance is computed using the mean decrease in Gini index and expressed relative to the maximum.